\chapter{Trampas conocidas}
\label{cap:trampas}

Cada entrada de este capítulo es un error que ya costó horas encontrar. Están agrupadas por
tecnología. La fuente viva es \archivo{docs/TRAMPAS.md}: cuando descubras una nueva, añade allí
una línea (síntoma, causa y arreglo) y, si merece explicación, un párrafo aquí.

\begin{idea}[Léelo antes de tocar]
Antes de cambiar algo de Qt, cxx-qt, mpv o YouTube, lee la sección que toque. Es la norma de
\archivo{docs/AGENTES.md} (§2): leer una trampa cuesta un minuto; reencontrarla, una tarde.
\end{idea}

\section{Compilación y cxx-qt}

\begin{description}[style=nextline, leftmargin=0pt]
  \item[Un archivo QML nuevo «no existe»] \emph{Síntoma:} QML dice que el tipo no existe.
  \emph{Causa:} no está en \cod{QML_FILES} (o en \cod{QML_SINGLETONS}) de
  \archivo{build.rs}. \emph{Arreglo:} añádelo. \cod{python3 scripts/docs.py check} lo comprueba.
  \item[Un puente Rust nuevo no aparece] Añádelo al \cod{.files([...])} de \archivo{build.rs}
  \emph{y} como \cod{mod} en \archivo{main.rs}.
  \item[Los nombres cambian en QML] \cod{#[auto_cxx_name]} convierte \cod{snake_case} en
  \cod{camelCase}: \cod{play_item} es \cod{Player.playItem} y \cod{vista_inicio} es
  \cod{AppSettings.vistaInicio}. Buscar el nombre de Rust desde QML no encuentra nada.
  \item[«Unsupported verbatim input in ForeignItem»] \emph{Causa:} metiste el \emph{cuerpo} de
  una función (o un método auxiliar) dentro del bloque \cod{extern "RustQt"}. \emph{Arreglo:}
  en el bloque va solo la firma de lo invocable; el cuerpo va en \cod{impl qobject::Nombre}.
\end{description}

\section{Qt Quick y QML}

\begin{description}[style=nextline, leftmargin=0pt]
  \item[\texttt{font.pixelSize} es entero] Un \cod{12.5} \emph{hace fallar la carga del
  componente entero}, no solo ese texto. Usa \cod{Math.round}.
  \item[No llames \texttt{palette} a tu propiedad] Choca con la propiedad \cod{palette} de
  \cod{Item} y QML ignora la tuya, sin avisar.
  \item[Un \texttt{id} tapa una propiedad del mismo nombre] En \archivo{NowPlayingView}, \cod{art}
  era una propiedad de la raíz \emph{y} el \cod{id} de la carátula: en los enlaces de la raíz,
  \cod{art} daba el elemento y la cuenta salía \cod{NaN}. Usa \cod{root.art}, o no repitas el
  nombre.
  \item[Cambiar un índice dentro del aviso de cambio de la misma propiedad] Hacerlo en
  \cod{onCurrentItemChanged} deja enlaces sin actualizar. Aplázalo: \cod{Qt.callLater(...)}.
  \item[Los clics atraviesan] Un \cod{TapHandler} normal solo toma un agarre \emph{pasivo}: el
  clic sigue hacia lo que haya debajo (cerrar «Sonando» pulsaba también el botón de opciones
  que estaba en el mismo sitio). En botones y opciones:
\begin{qml}
TapHandler { gesturePolicy: TapHandler.ReleaseWithinBounds; onTapped: ... }
\end{qml}
  Un \cod{MouseArea} debajo no lo arregla.
  \item[La rueda atraviesa las capas de encima] Con Opciones o el menú abiertos, la rueda
  desplaza la sección de debajo: un \cod{Flickable} que llega al tope, o que no se desplaza,
  deja pasar la rueda. Solución: bajo la capa, un \cod{MouseArea} con
  \cod{acceptedButtons: Qt.NoButton} y \cod{onWheel: wheel => wheel.accepted = true}.
  \item[\texttt{Window.window} dentro de un \texttt{TapHandler} es \texttt{null}] Los
  \emph{handlers} no son \cod{Item}. Usa el del elemento:
  \cod{delegado.Window.window.openMenu(...)}.
  \item[Arrastrar dentro de una lista la desplaza] En vez de mover la fila. Usa un
  \cod{MouseArea} con \cod{preventStealing: true} (como \archivo{QueueList}).
  \item[«Hover» sobre algo que se mueve] Un \cod{HoverHandler}, un \cod{MouseArea} con
  \cod{cursorShape} o un \cod{TapHandler} \emph{por elemento} sobre contenido animado obligan a
  repasar esa parte de la escena en cada fotograma. En la letra llevaba la CPU de ~5\,\% a
  ~25\,\%. Un único gestor para toda la vista, con \cod{childAt}, y sin «hover».
  \item[Rotaciones 3D] Un ángulo positivo en el eje Y \emph{aleja el lado derecho}. Y las
  rotaciones grandes (más de ~70°) acercan la pieza a la cámara y la agrandan: para abrir la
  caja se usa un pliegue plano (\cod{Scale} en X).
  \item[Máscaras de \texttt{MultiEffect}] Sin \cod{maskThresholdMin: 0.5} y
  \cod{maskSpreadAtMin: 1} la máscara recorta en seco en vez de seguir el degradado.
  \item[\texttt{TintProbe} (un \texttt{Canvas})] Cada lienzo debe llamar a \cod{loadImage}
  aunque otro ya tenga la imagen, y repintar con su propio \cod{Timer}; si no, la segunda
  sonda con la misma carátula no pinta. Una sonda recortada (fuera de la zona visible) tampoco
  pinta: usa la de \archivo{Main.qml} (\cod{playingProbe}).
  \item[Una \texttt{Shape} con \texttt{CurveRenderer} no se redibuja tras ocultar la ventana]
  Al liberar los gráficos (\archivo{window.cpp}) se pierde. \archivo{Icon.qml} se recrea con un
  \cod{Loader} cada vez que sube \cod{Theme.renderEpoch}; una \cod{Shape} nueva con ese
  renderizador necesita lo mismo.
  \item[\texttt{ShapePath.fillItem} no usa el tamaño del elemento] Coloca la textura en la
  esquina de la forma, con su tamaño en \emph{píxeles físicos} (una imagen de 544 px en Retina
  ocupa 272). En una forma mayor, el último píxel se repite por la derecha y por abajo: la
  carátula salía «corrida». Escálala con \cod{fillTransform}, como \archivo{Cover.qml}. La
  captura sin ventana (\cod{offscreen}) no pinta el relleno: compruébalo con ventana.
  \item[Una \texttt{ListView} vertical pone cada delegado en \texttt{x = 0}] Aunque el delegado
  diga otra \cod{x}. El margen va en \cod{leftMargin} de la lista (el Muro salía pegado al
  borde).
  \item[Letra en japonés o chino] Partir la línea por espacios deja una sola «palabra»
  enorme que se encoge en vez de pasar al renglón siguiente, y unir las sílabas con espacios
  separa los caracteres. Cada sílaba lleva \cod{join} y los renglones se parten entre
  caracteres CJK (\cod{rows} de \archivo{LyricsView}).
\end{description}

\begin{trampa}[La más cara: el \emph{hover} en la letra]
Es la trampa que más CPU costó, porque no falla: todo funciona y solo el consumo lo delata. Si
algo que se anima va lento o gasta de más, \emph{lo primero que hay que mirar es si hay un
gestor de puntero por elemento encima}.
\end{trampa}

\section{YouTube, yt-dlp y mpv}

\begin{description}[style=nextline, leftmargin=0pt]
  \item[No castigues a YouTube] Muchas peticiones seguidas y rechaza la IP un rato:
  «Requested format is not available», 403, «Sign in to confirm you're not a bot». Espera unos
  minutos y compruébalo con \cod{yt-dlp} directamente antes de culpar al código. Agrupa las
  pruebas que hacen peticiones reales.
  \item[mpv y las URL de googlevideo] Las de yt-dlp dan 403 con el cliente HTTP de mpv y se
  cargan como \cod{lavf://URL} (ffmpeg); las de la vía rápida se atascan con ffmpeg y se
  cargan tal cual. \cod{resolver::run_ytdlp} deja el prefijo según el origen.
  \item[La vía rápida necesita «visitor data»] Sin él YouTube responde
  \cod{LOGIN_REQUIRED} (\archivo{stream.rs}).
  \item[yt-dlp sin sesión] Con los clientes por defecto YouTube pide verificación anti-bot en
  pistas de YouTube Music. Se usan \cod{web_embedded,mweb} (\cod{resolver::ytdlp_clients}).
  \item[yt-dlp con sesión] Solo \cod{web_music} da las pistas reservadas a cuentas y los
  formatos Premium; con \cod{player_skip=webpage,configs} tarda ~1,7 s en vez de ~3 s.
  Diagnóstico: \cod{--probe playable ID} y \cod{--probe account}.
  \item[yt-dlp reescribe el archivo de \texttt{--cookies}] Al terminar. \emph{Nunca} le pases el
  \cod{cookies.txt} original: usa \cod{account::cookies_copy}.
  \item[mpv solo avisa de la pausa cuando cambia] Al empezar una pista hay que aplicar la
  última conocida (\cod{mpv_paused}); si no, «sonando» se queda desfasado.
  \item[El historial de YouTube Music] Los avisos de escucha necesitan \cod{c=WEB_REMIX}; sin
  él, la canción aparece en el historial de YouTube y no en el de YouTube Music.
  \item[La cola \texttt{next} de un álbum] Trae versiones en vídeo, se corta en ~50 y mezcla
  canciones ajenas: se usa la página \cod{VL<id>} (ADR-06).
\end{description}

\section{macOS y Linux}

\begin{description}[style=nextline, leftmargin=0pt]
  \item[Minimizar no libera memoria gráfica en macOS] Ocultar la ventana sí
  (\archivo{window.cpp}).
  \item[Las apps gráficas de macOS no heredan el \texttt{PATH} de la terminal] Por eso
  \cod{audio::search_path} añade \cod{/opt/homebrew/bin} y \cod{/usr/local/bin}: mpv y yt-dlp
  no se encontrarían.
  \item[El proxy del sistema] Qt ignora el proxy de macOS (\archivo{window.cpp}); con un WPAD
  roto cada descarga fallaba con «NSURLErrorDomain -1003». Se respetan \cod{HTTPS_PROXY} y
  \cod{HTTP_PROXY}.
  \item[Inicio de sesión en WebKit] La prevención de rastreo pierde las cookies en las
  redirecciones de Google y YouTube acaba en \cod{/oops}; se desactiva.
  \item[Disfrazar el agente de usuario de Chrome en Qt WebEngine] Google lo bloquea («este
  navegador puede no ser seguro») porque contradice \cod{Sec-CH-UA}. No lo toques.
  \item[Los sockets Unix admiten ~104 bytes de ruta] En pruebas de Discord usa una ruta
  relativa corta (\cod{audio.rs} cae a \cod{/tmp} por la misma razón).
\end{description}

\section{Medir}

\begin{description}[style=nextline, leftmargin=0pt]
  \item[Mide siempre en release] Depuración cambia la memoria y los tiempos.
  \item[Mide la CPU con tiempo acumulado] Con \cod{ps -o time=} al principio y al final de
  25 s, y comprueba en la línea \cod{estado:} que sonaba la canción esperada (Inicio cambia de
  una ejecución a otra). \cod{ps %cpu} es una foto de un instante y engaña.
  \item[Memoria con \texttt{footprint}] La que muestra el Monitor de Actividad. Cualquier
  ventana de Qt Quick que se redibuje reserva ~100 MB en el controlador gráfico.
\end{description}

\section{Trampas del trabajo con el repositorio}

Las normas de \archivo{docs/AGENTES.md} existen por errores de método, no de código:

\begin{itemize}
  \item \textbf{Los datos del propietario no se tocan.} Una prueba que lee o borra
  \cod{<config>/musicplayer/} se carga la sesión de la persona. Usa un \cod{HOME} temporal.
  \item \textbf{Nunca importes las cookies del navegador} en una prueba: se prueba el camino
  con sesión con un \cod{cookies.txt} falso.
  \item \textbf{\texttt{MUSICPLAYER\_SILENT=1} siempre} en las pruebas automáticas.
  \item \textbf{Una rama por tarea}: no se trabaja en \cod{master}.
  \item \textbf{No presentes como probado lo que no lo está}: Linux, Wayland, el sonido real,
  con sesión y la app empaquetada se dicen explícitamente cuando no se verificaron.
\end{itemize}

\begin{ejercicio}
Elige tres trampas de este capítulo y, para cada una, localiza en el código el sitio donde se
evita (por ejemplo, \cod{preventStealing: true} en \archivo{QueueList.qml}). Si quitaras esa
línea, ¿cuál sería el síntoma exacto?
\end{ejercicio}
